\abstract{РЕФЕРАТ}

Объем работы равен \formbytotal{lastpage}{страниц}{е}{ам}{ам}. Работа содержит \formbytotal{figurecnt}{иллюстраци}{ю}{и}{й}, \formbytotal{tablecnt}{таблиц}{у}{ы}{}, \arabic{bibcount} библиографических источников и \formbytotal{числоПлакатов}{лист}{}{а}{ов} графического материала. Количество приложений – 2. Графический материал представлен в приложении А. Фрагменты исходного кода представлены в приложении Б.

Перечень ключевых слов: файловая система, виртуальная файловая система, FAT32, ISO, CDFS, дескриптор, блок, диск, сектор, файл, каталог, классы, таблица путей, таблица, цепочка, загрузочный сектор, имя файла, расширение файла, данные, байты, биты, запись, массив, Big-Endian, Little-Endian, экстент, системная область, парсер, исключения, ошибка, методы, фрагментация, наследование, безопасность, журналирование.

Объектом разработки является виртуальная файловая система, поддерживающая файловые системы FAT32 и ISO9660(CDFS).

Целью выпускной квалификационной работы является создание виртуальной файловой системы, а также изучение структуры FAT32 и CDFS.

В процессе создания виртуальной файловой системы были выделены основные операции над файлами, на основе которых был спроектирован интерфейс виртуальной файловой системы, созданы классы и методы конкретных файловых систем, реализующие интерфейс виртуальной файловой системы. Реализован командный интерпретатор, обеспечивающий интерактивное взаимодействие с виртуальной файловой системой.

При разработке виртуальной файловой системы использовался язык программирования "<Common Lisp">.

Разработанная виртуальная файловая система была успешно внедрена в "<os-pi">.

\selectlanguage{english}
\abstract{ABSTRACT}
  
The volume of work is \formbytotal{lastpage}{page}{}{s}{s}. The work contains \formbytotal{figurecnt}{illustration}{}{s}{s}, \formbytotal{tablecnt}{table}{}{s}{s}, \arabic{bibcount} bibliographic sources and \formbytotal{числоПлакатов}{sheet}{}{s}{s} of graphic material. The number of applications is 2. The graphic material is presented in annex A. Fragments of the source code are presented in annex B.

List of keywords: file system, virtual file system, FAT32, ISO, CDFS, descriptor, block, disk, sector, file, directory, classes, path table, table, chain, boot sector, file name, file extension, data, bytes, bits, record, array, Big-Endian, Little-Endian, extent, system area, parser, exceptions, error, methods, fragmentation, inheritance, security, journaling.

The object of the development is a virtual file system supporting the FAT32 and ISO9660 (CDFS) file systems.

The aim of the graduation thesis is the creation of a virtual file system, as well as studying the structure of FAT32 and CDFS.

In the process of creating the virtual file system, the main file operations were identified, based on which the virtual file system interface was designed, and classes and methods of specific file systems implementing the virtual file system interface were created. A command interpreter was implemented, providing interactive interaction with the virtual file system.

The programming language <<Common Lisp>> was used in the development of the virtual file system.

The developed virtual file system was successfully implemented in <<os-pi>>.
\selectlanguage{russian}
